% SOURCE_AUTHORITY: sga5_fr_workpass.tex, lines 10179--10234 (LF-normalized unit).
% SOURCE_SHA256: 791F4EFFC5E02832D5D77ED03518C8156D6F07E4C8238B03545DB93D883FBB28
Vamos agora dar uma variante relativa de (5.4) e (5.5). Com as notações anteriores, seja $u:S\to T$ um morfismo de esquemas. Ele se põe
\[
A^\bullet(X/T)=\bigoplus_i \R^{2i}(u\circ f)_*(A_X)
\]
e, igualmente,
\[
A^\bullet(S/T)=\bigoplus_i \R^{2i}u_*(A_S).
\]
Finalmente, para tudo $A_S$-módulo $L$, é colocado
\[
Q(S/T,L)=\bigoplus_{n,i}\R^n u_*(L\otimes\mu_S^{\otimes i})
\]
e
\[
Q(X/T,f^*L)=\bigoplus_{n,i}\R^n(u\circ f)_*(f^*L\otimes\mu_X^{\otimes i}).
\]

Para todo par $(i,j)$ de inteiros, com $1\le i\le m$, $1\le j\le p_i$, a imagem da classe de Chern $c_j(E_i)$ pelo canônico morfismo
\[
H^{2j}(X,\mu_X^{\otimes j})
\to H^0(T,\R^{2j}(u\circ f)_*(\mu_X^{\otimes j}))
\]
define um $A_T$-morfismo
\[
c_{ij}:A_T\to \R^{2j}(u\circ f)_*(\mu_X^{\otimes j}).
\]
A data de $c_{ij}$ e a imagem inversa $f^*$ permitem definir um morfismo de $A^*(S/T)$-álgebra
\[
\theta_1:A^\bullet(S/T)[T_{i,j_i}]_{(i,j_i)}\to A^\bullet(X/T).
\]
Denotando por $c_i$ as seções de $A^*(S/T)$ definidas pelas classes de Chern $c_i(E)$, se vê, como em (5.4), que $\theta_1$ é anulada sobre o ideal $J$ gerado localmente pelas relações
\[
\sum_{i_1+\cdots+i_m=p}T_{1,i_1}\cdots T_{m,i_m}=c_p\qquad(p\in\mathbb N).
\]
Daí um morfismo de $A^*(S/T)$-álgebra
\[
\theta:A^\bullet(S/T)[T_{i,j_i}]_{(i,j_i)}/J\to A^\bullet(X/T).
\]

\begin{statement}{Proposição 5.6}
Com as notações anteriores:
\begin{enumerate}[label=\alph*)]
\item O morfismo de $A^*(S/T)$-álgebra
\[
\theta:A^\bullet(S/T)[T_{i,j_i}]_{(i,j_i)}/J\to A^\bullet(X/T)
\]
é um isomorfismo.
\item Para todo $A_S$-módulo $L$, o canônico morfismo de $A^*(X/T)$-módulos
\[
A^\bullet(X/T)\otimes_{A^\bullet(S/T)}Q(S/T,L)\to Q(X/T,f^*L)
\]
é um isomorfismo.
\end{enumerate}
\end{statement}

\emph{Demonstração.} Basta fazer (5.4) e (5.5), respectivamente.
